\documentclass[11pt,a4paper]{article}
\usepackage{fontspec,amsmath,amssymb,amsthm,microtype}
\setmainfont{lmroman10-regular.otf}[BoldFont=lmroman10-bold.otf,ItalicFont=lmroman10-italic.otf,BoldItalicFont=lmroman10-bolditalic.otf]
\usepackage{xcolor}
\usepackage[most]{tcolorbox}
\usepackage{mathrsfs}
\providecommand{\lt}{<}
\providecommand{\gt}{>}
\setmonofont{DejaVuSansMono.ttf}[BoldFont=DejaVuSansMono-Bold.ttf,ItalicFont=DejaVuSansMono-Oblique.ttf,BoldItalicFont=DejaVuSansMono-BoldOblique.ttf,Scale=0.86]
\newfontfamily\leanfallback{STIXTwoMath-Regular.otf}[Scale=0.86]
\definecolor{leanrule}{gray}{0.62}
\definecolor{leanback}{gray}{0.962}
\definecolor{leantitle}{gray}{0.22}
\definecolor{leanmuted}{gray}{0.38}
\newtcolorbox{leanbox}[1]{enhanced,breakable,colback=leanback,colframe=leanrule,
  boxrule=0.4pt,arc=0.6mm,left=2.4mm,right=2.4mm,top=1.8mm,bottom=1.8mm,
  before skip=0.8em,after skip=1em,fonttitle=\small\bfseries,coltitle=leantitle,
  colbacktitle=leanback,titlerule=0.3pt,title={#1}}
\newlength{\leancharwidth}
\newenvironment{leancode}{\par\smallskip\begingroup\ttfamily\footnotesize\raggedright
  \settowidth{\leancharwidth}{M}\setlength{\parskip}{0pt}\setlength{\parindent}{0pt}
  \frenchspacing}{\par\endgroup\smallskip}
\newcommand{\leanline}[2]{\hangindent=\dimexpr#1\leancharwidth+4\leancharwidth\relax
  \hangafter=1\noindent\hspace*{#1\leancharwidth}#2\par}
\newcommand{\leansource}[1]{\par{\footnotesize\color{leanmuted}#1}\par}
\newcommand{\leanlabel}[1]{\par\smallskip\noindent{\small\bfseries #1}\par}
\usepackage[margin=25mm]{geometry}
\usepackage[unicode,hidelinks]{hyperref}
\hypersetup{pdfauthor={LangevinLean Project}}
\setlength{\emergencystretch}{3em}
\setlength{\parskip}{0.4em}
\setlength{\parindent}{0pt}
\allowdisplaybreaks[2]
\raggedbottom
\title{\texorpdfstring{Moments and coupling error of the actual random MALA path}{Moments and coupling error of the actual random MALA path}}
\author{LangevinLean Project}
\date{}
\begin{document}\maketitle
Draft mathematical exposition; these notes may contain mistakes.


The canonical sampled MALA recursion is converted to a genuine \(L^2\) random variable. Its marginal remains exactly Gibbs, and its mean distance from the synchronous Euler coordinate satisfies the proved cumulative rejection estimate.


\section{\texorpdfstring{Actual representatives, Gibbs laws and square integrability}{Actual representatives, Gibbs laws and square integrability}}


\textbf{Claim.} The sampled input and paired states are measurable. The first coordinate is the actual Euler recursion; the second has Gibbs law and finite second moment, hence defines an actual \(L^2\) variable with that same law.



\textbf{Proof.} Compose the measurable sampled update through finite recursion. Projecting its exact law gives the established MALA marginal, which is Gibbs by invariance at each finite step. Pull the Gibbs finite second moment through this exact pushforward to obtain square integrability; passing to the \(L^2\) class preserves its almost-everywhere representative and law.

\begin{leanbox}{Lean: \texttt{input\_\allowbreak{}aemeasurable}, \texttt{state\_\allowbreak{}aemeasurable}, \texttt{euler} and 6 more}
\leansource{\texttt{input\_\allowbreak{}aemeasurable} \quad WassersteinMixing/\allowbreak{}SemigroupRandomMALABounds.lean, lines 21–25}
\begin{leancode}
\leanline{0}{theorem\ input\_aemeasurable\ (r\ :\ ℝ)\ :}
\leanline{4}{AEMeasurable\ (input\ P\ V\ hV\ r)\ (probability\ P\ V)}
\end{leancode}
\leansource{\texttt{state\_\allowbreak{}aemeasurable} \quad WassersteinMixing/\allowbreak{}SemigroupRandomMALABounds.lean, lines 27–33}
\begin{leancode}
\leanline{0}{theorem\ state\_aemeasurable\ (r\ :\ ℝ)\ (n\ :\ ℕ)\ :}
\leanline{4}{AEMeasurable\ (state\ P\ V\ hV\ r\ n)\ (probability\ P\ V)}
\end{leancode}
\leansource{\texttt{euler} \quad WassersteinMixing/\allowbreak{}SemigroupRandomMALABounds.lean, lines 36–39}
\begin{leancode}
\leanline{0}{def\ euler\ (r\ :\ ℝ)\ (n\ :\ ℕ)\ :\ Lp\ (E\ P.dimension)\ 2\ (actualProbability\ P\ V)\ :=\ by}
\leanline{2}{let\ :=\ actualProbability\_isProbability\ P\ V\ hV}
\leanline{2}{exact\ ULA.Hilbert.lpChain\ (actualProbability\ P\ V)\ V\ r\ hV.gradientLipschitz}
\leanline{4}{(actualInitial\ P\ V\ hV)\ (gaussian\ P\ V\ hV\ r)\ n}
\end{leancode}
\leansource{\texttt{state\_\allowbreak{}fst\_\allowbreak{}eq\_\allowbreak{}euler} \quad WassersteinMixing/\allowbreak{}SemigroupRandomMALABounds.lean, lines 41–55}
\begin{leancode}
\leanline{0}{theorem\ state\_fst\_eq\_euler\ (r\ :\ ℝ)\ (n\ :\ ℕ)\ :}
\leanline{4}{(fun\ ω\ ↦\ (state\ P\ V\ hV\ r\ n\ ω).1)\ =ᵐ[probability\ P\ V]}
\leanline{6}{(fun\ ω\ ↦\ euler\ P\ V\ hV\ r\ n\ ω.1)}
\end{leancode}
\leansource{\texttt{state\_\allowbreak{}snd\_\allowbreak{}law} \quad WassersteinMixing/\allowbreak{}SemigroupRandomMALABounds.lean, lines 58–67}
\begin{leancode}
\leanline{0}{theorem\ state\_snd\_law\ (r\ :\ ℝ)\ (hr\ :\ 0\ <\ r)\ (n\ :\ ℕ)\ :}
\leanline{4}{(probability\ P\ V).map\ (fun\ ω\ ↦\ (state\ P\ V\ hV\ r\ n\ ω).2)\ =\ ULA.gibbsMeasure\ V}
\end{leancode}
\leansource{\texttt{state\_\allowbreak{}snd\_\allowbreak{}memLp} \quad WassersteinMixing/\allowbreak{}SemigroupRandomMALABounds.lean, lines 69–75}
\begin{leancode}
\leanline{0}{theorem\ state\_snd\_memLp\ (r\ :\ ℝ)\ (hr\ :\ 0\ <\ r)\ (n\ :\ ℕ)\ :}
\leanline{4}{MemLp\ (fun\ ω\ ↦\ (state\ P\ V\ hV\ r\ n\ ω).2)\ 2\ (probability\ P\ V)}
\end{leancode}
\leansource{\texttt{mala} \quad WassersteinMixing/\allowbreak{}SemigroupRandomMALABounds.lean, lines 78–79}
\begin{leancode}
\leanline{0}{def\ mala\ (r\ :\ ℝ)\ (hr\ :\ 0\ <\ r)\ (n\ :\ ℕ)\ :\ Lp\ (E\ P.dimension)\ 2\ (probability\ P\ V)\ :=}
\leanline{2}{(state\_snd\_memLp\ P\ V\ hV\ r\ hr\ n).toLp\ (fun\ ω\ ↦\ (state\ P\ V\ hV\ r\ n\ ω).2)}
\end{leancode}
\leansource{\texttt{mala\_\allowbreak{}coe} \quad WassersteinMixing/\allowbreak{}SemigroupRandomMALABounds.lean, lines 81–83}
\begin{leancode}
\leanline{0}{theorem\ mala\_coe\ (r\ :\ ℝ)\ (hr\ :\ 0\ <\ r)\ (n\ :\ ℕ)\ :}
\leanline{4}{mala\ P\ V\ hV\ r\ hr\ n\ =ᵐ[probability\ P\ V]\ fun\ ω\ ↦\ (state\ P\ V\ hV\ r\ n\ ω).2}
\end{leancode}
\leansource{\texttt{mala\_\allowbreak{}law} \quad WassersteinMixing/\allowbreak{}SemigroupRandomMALABounds.lean, lines 85–87}
\begin{leancode}
\leanline{0}{theorem\ mala\_law\ (r\ :\ ℝ)\ (hr\ :\ 0\ <\ r)\ (n\ :\ ℕ)\ :}
\leanline{4}{(probability\ P\ V).map\ (mala\ P\ V\ hV\ r\ hr\ n)\ =\ ULA.gibbsMeasure\ V}
\end{leancode}
\leanlabel{How the Lean proof follows the argument}
The initial measurability lemmas justify the sampled pushforwards. \texttt{state\_\allowbreak{}snd\_\allowbreak{}law} and \texttt{state\_\allowbreak{}snd\_\allowbreak{}mem\allowbreak{}Lp} establish the true marginal and moment before defining \texttt{mala}. \texttt{mala\_\allowbreak{}coe} and \texttt{mala\_\allowbreak{}law} retain the representative and exact law.
\end{leanbox}


Earlier results used here: \cite[Section~3.19.1]{LangevinLean2026ULASC}; \cite[Section~3.20.1]{LangevinLean2026ULASC}; \cite[Section~3.20.4]{LangevinLean2026ULASC}; \cite[Section~3.23.1]{LangevinLean2026ULASC}; \cite[Section~3.26.3]{LangevinLean2026ULASC}; \cite[Section~3.36.17]{LangevinLean2026ULASC}; \cite[Section~3.39.6]{LangevinLean2026ULASC}; \cite[Section~3.42.1]{LangevinLean2026ULASC}.

\section{\texorpdfstring{The actual accumulated rejection distance}{The actual accumulated rejection distance}}


\textbf{Claim.} Under the stated small-step conditions, the expected distance between the actual MALA state and the lifted Euler state is at most \(nCr\sqrt r\), where \(C\) is the proved Gaussian rejection constant for the potential.



\textbf{Proof.} Both coordinates are square integrable, so their difference norm is integrable. The actual joint law equals the established coupled-chain law. Transfer its cumulative rejection-cost estimate through that exact law and through the almost-everywhere representatives of the \(L^2\) states.

\begin{leanbox}{Lean: \texttt{distance\_\allowbreak{}integrable}, \texttt{distance\_\allowbreak{}integral\_\allowbreak{}le}}
\leansource{\texttt{distance\_\allowbreak{}integrable} \quad WassersteinMixing/\allowbreak{}SemigroupRandomMALABounds.lean, lines 89–99}
\begin{leancode}
\leanline{0}{theorem\ distance\_integrable\ (r\ :\ ℝ)\ (hr\ :\ 0\ <\ r)\ (hr1\ :\ r\ ≤\ 1)}
\leanline{4}{(hsmall\ :\ r\ ≤\ 2\ /\ (P.alpha\ +\ P.beta))\ (n\ :\ ℕ)\ :}
\leanline{4}{Integrable\ (fun\ ω\ ↦\ ‖(state\ P\ V\ hV\ r\ n\ ω).1\ -\ (state\ P\ V\ hV\ r\ n\ ω).2‖)}
\leanline{6}{(probability\ P\ V)}
\end{leancode}
\leansource{\texttt{distance\_\allowbreak{}integral\_\allowbreak{}le} \quad WassersteinMixing/\allowbreak{}SemigroupRandomMALABounds.lean, lines 103–114}
\begin{leancode}
\leanline{0}{theorem\ distance\_integral\_le\ (r\ :\ ℝ)\ (hr\ :\ 0\ <\ r)\ (hr1\ :\ r\ ≤\ 1)}
\leanline{4}{(hsmall\ :\ r\ ≤\ 2\ /\ (P.alpha\ +\ P.beta))\ (n\ :\ ℕ)\ :}
\leanline{4}{(∫\ ω,\ ‖(state\ P\ V\ hV\ r\ n\ ω).1\ -\ (state\ P\ V\ hV\ r\ n\ ω).2‖\ ∂probability\ P\ V)\ ≤}
\leanline{6}{(n\ :\ ℝ)\ *\ ULA.malaRejectionConstant\ V\ ⟨P.beta,\ P.beta\_pos.le⟩\ *\ r\ *\ Real.sqrt\ r}
\end{leancode}
\leanlabel{How the Lean proof follows the argument}
\texttt{distance\_\allowbreak{}integrable} uses the two true \(L^2\) coordinates. \texttt{distance\_\allowbreak{}integral\_\allowbreak{}le} rewrites the actual expectation under \texttt{state\_\allowbreak{}law} and invokes the proved coupled MALA/\allowbreak{}Euler rejection estimate.
\end{leanbox}


Earlier results used here: \cite[Section~3.34.22]{LangevinLean2026ULASC}; \cite[Section~3.36.25]{LangevinLean2026ULASC}; \cite[Section~3.36.26]{LangevinLean2026ULASC}; \cite[Section~3.39.6]{LangevinLean2026ULASC}.

\begin{thebibliography}{ULA-LSI}
\bibitem[ULA-SC]{LangevinLean2026ULASC}
LangevinLean Project. \emph{Unadjusted Langevin under Strong Convexity}. LL-01, 2026. Draft report, content revision d34ead06987c. \url{https://odakyildiz.github.io/LangevinLean/docs/pdf/UnadjustedLangevin/FullProof.pdf}.
\end{thebibliography}

\end{document}
